\section{生成对抗网络（GAN）}

\subsection{GAN 的基本思想}

\begin{figure}[H]
\centering
\includegraphics[width=0.95\textwidth]{slides-images/slide-048.jpg}
\caption{GAN 的基本架构：Generator 将噪声 $Z$ 转化为数据 $X$，Discriminator 判断数据是真实还是生成的\protect\footnotemark}
\end{figure}
\footnotetext{Slides 第49页。}

GAN 属于\textbf{隐式密度模型}——放弃显式建模 $p(X)$，转而通过一个博弈过程隐式地学习数据分布。

GAN 包含两个神经网络：
\begin{itemize}
    \item \textbf{Generator $G$}：输入噪声 $Z \sim \mathcal{N}(0,I)$，输出生成数据 $G(Z)$
    \item \textbf{Discriminator $D$}：输入数据（真实或生成），输出「真实」的概率 $D(X) \in [0,1]$
\end{itemize}

\subsection{目标函数}

\begin{figure}[H]
\centering
\includegraphics[width=0.95\textwidth]{slides-images/slide-051.jpg}
\caption{GAN 的 minimax 目标函数\protect\footnotemark}
\end{figure}
\footnotetext{Slides 第52页。}

GAN 的训练是一个 minimax 博弈：

$$\min_G \max_D \left[ \mathbb{E}_{X \sim p_{\text{data}}}[\log D(X)] + \mathbb{E}_{Z \sim p(Z)}[\log(1 - D(G(Z)))] \right]$$

\begin{itemize}
    \item \textbf{Discriminator 的目标}（最大化）：正确分类真实数据为真（$D(X) \to 1$），生成数据为假（$D(G(Z)) \to 0$）
    \item \textbf{Generator 的目标}（最小化）：让 Discriminator 把生成数据误判为真（$D(G(Z)) \to 1$）
\end{itemize}

\subsection{训练过程}

\begin{figure}[H]
\centering
\includegraphics[width=0.95\textwidth]{slides-images/slide-053.jpg}
\caption{GAN 的交替训练：先更新 Discriminator 若干步，再更新 Generator 一步\protect\footnotemark}
\end{figure}
\footnotetext{Slides 第54页。}

训练交替进行两个步骤：
\begin{enumerate}
    \item \textbf{训练 Discriminator}：固定 Generator，用真实数据和生成数据训练 Discriminator 做二分类
    \item \textbf{训练 Generator}：固定 Discriminator，通过 Discriminator 的反馈优化 Generator
\end{enumerate}

\begin{warningbox}{GAN 训练的梯度消失问题}
在训练初期，Generator 产生的样本质量极差，Discriminator 很容易将其识别为假。此时 $D(G(Z)) \approx 0$，而 $\log(1 - D(G(Z)))$ 在 $D(G(Z)) \approx 0$ 处梯度非常小。这意味着 Generator 在训练初期几乎收不到有效的梯度信号。实践中常用的替代目标是最大化 $\log D(G(Z))$（而非最小化 $\log(1 - D(G(Z)))$），因为前者在 $D(G(Z)) \approx 0$ 处梯度更大。
\end{warningbox}

\begin{figure}[H]
\centering
\includegraphics[width=0.95\textwidth]{slides-images/slide-055.jpg}
\caption{Generator 的两种目标函数对比：原始目标在训练初期梯度平坦，替代目标梯度更强\protect\footnotemark}
\end{figure}
\footnotetext{Slides 第56页。}

\subsection{GAN 的训练难点}

\begin{figure}[H]
\centering
\includegraphics[width=0.95\textwidth]{slides-images/slide-058.jpg}
\caption{GAN 训练的典型曲线：Generator 和 Discriminator 的损失曲线难以解读，无法判断训练进展\protect\footnotemark}
\end{figure}
\footnotetext{Slides 第59页。}

\begin{warningbox}{GAN 训练是出了名的不稳定}
\begin{itemize}
    \item \textbf{没有有意义的损失曲线}：与其他模型不同，GAN 的损失值不能直观反映生成质量。数百篇论文试图解决这个问题，但至今仍未有完美方案。
    \item \textbf{模式坍缩}（Mode Collapse）：Generator 可能只学会生成少数几种样本，忽略数据分布的多样性。
    \item \textbf{训练不收敛}：Generator 和 Discriminator 的对抗可能导致训练震荡而非收敛。
    \item \textbf{超参数敏感}：对学习率、网络架构、训练步数等超参数极为敏感。
\end{itemize}
\end{warningbox}

\subsection{GAN 生成的图像质量}

尽管训练困难，GAN 在生成图像质量上曾长期领先。

\begin{figure}[H]
\centering
\includegraphics[width=0.95\textwidth]{slides-images/slide-060.jpg}
\caption{GAN 生成的高质量人脸图像（StyleGAN 系列）\protect\footnotemark}
\end{figure}
\footnotetext{Slides 第61页。}

\begin{figure}[H]
\centering
\includegraphics[width=0.95\textwidth]{slides-images/slide-062.jpg}
\caption{BigGAN 生成的高分辨率类别条件图像：从犬类到食物的多种类别\protect\footnotemark}
\end{figure}
\footnotetext{Slides 第63页。}

\begin{figure}[H]
\centering
\includegraphics[width=0.95\textwidth]{slides-images/slide-067.jpg}
\caption{GAN 的 latent space 插值：在两个 latent code 之间线性插值产生平滑过渡\protect\footnotemark}
\end{figure}
\footnotetext{Slides 第68页。}

GAN 在约 2016--2021 年间是图像生成的主流方法。它的优势在于：
\begin{itemize}
    \item 单次前向传播即可生成样本（推理速度快）
    \item 生成的样本清晰、锐利
\end{itemize}

\subsection{条件 GAN}

通过向 Generator 和 Discriminator 输入条件信息 $Y$（如类别标签、文本描述），可以将 GAN 扩展为条件生成模型。

\begin{figure}[H]
\centering
\includegraphics[width=0.95\textwidth]{slides-images/slide-070.jpg}
\caption{条件 GAN 的应用：语义分割图到逼真图像的转换（Pix2Pix）\protect\footnotemark}
\end{figure}
\footnotetext{Slides 第71页。}

\subsection{GAN vs VAE 的权衡}

\begin{figure}[H]
\centering
\includegraphics[width=0.95\textwidth]{slides-images/slide-073.jpg}
\caption{VAE vs GAN：VAE 有 latent vectors 和密度估计但样本模糊；GAN 样本清晰但训练困难\protect\footnotemark}
\end{figure}
\footnotetext{Slides 第74页。}

\begin{importantbox}{VAE 和 GAN 的权衡}
\begin{itemize}
    \item \textbf{VAE}：可以计算（近似的）$p(X)$；有结构化的 latent space（可做编码和插值）；但样本通常偏模糊
    \item \textbf{GAN}：样本质量高、清晰锐利；推理速度快（单次前向）；但训练困难、无有效损失指标、容易模式坍缩、没有 encoder（无法从 $X$ 到 $Z$ 的映射）
\end{itemize}
\end{importantbox}

\subsection{本章小结}

GAN 通过 Generator 和 Discriminator 的对抗训练隐式地学习数据分布。它放弃了密度估计的能力，换来了更清晰的生成样本。GAN 在 2016--2021 年间主导了图像生成领域，但其训练不稳定性是重大缺陷。后来被扩散模型取代成为主流方法。

%% ========== 总结与延伸 ==========

